\section{Online observations and ordinary marginals}
\label{sec:marginals}

We track the conditional location of one hidden card while revealing
the paths of earlier cards in a list. The quadratic error and the
successive comparison of conditional location laws follow the
swap-or-not approach of Hoang, Morris, and Rogaway
\cite[Section~3, Theorem~3 and Lemma~4]{hmr2014}. Here the next
direction is constrained by its recent predecessors, so we state
the information used in the calculation explicitly. The first
application below produces one partition whose complementary
marginals are all close to uniform. Later applications use the same
calculation with other sources of balance.

\subsection{A chronological law and a hidden card}

Fix an integer $T\ge0$. Let $\mathbf v=(v_1,\ldots,v_T)$ be a random schedule of nonzero
directions. Set $Z_0=\id$, and suppose that, conditional on the
complete schedule, the increments of $Z$ are independent uniform
switches in directions $v_1,\ldots,v_T$. This law can be generated
in time order: sample the next direction from its conditional law
given the preceding directions, and then use fresh fair coins on
its matching.

Let $\Xi$ be initial data independent of the joint process
$(\mathbf v,(Z_t)_{t=0}^T)$, and let
$C=(c_1,\ldots,c_q)$ be a list of distinct labels determined by
$\Xi$, with fixed length $0\le q\le n$. A fixed list corresponds to
deterministic $\Xi$; a sampled partition may be included in $\Xi$
when it determines the list. For its $j$th label define the exact
online field
\begin{equation}\label{eq:marginal-filtration}
 \mathcal F_t^{j}
 =\sigma\bigl(\Xi,v_1,\ldots,v_t,\,
        Z_s(c_i):0\le s\le t,\ i<j\bigr).
\end{equation}
Let $D_t^j$ be the positions not occupied by the observed labels
$c_1,\ldots,c_{j-1}$ at time $t$, and put $r_j=n-j+1$. Define
\begin{equation}\label{eq:marginal-energy}
 p_t^j(x)=\Pr\{Z_t(c_j)=x\mid\mathcal F_t^j\},\qquad
 \Delta_t^j=\sum_{x\in D_t^j}\left(p_t^j(x)-\frac1{r_j}\right)^2.
\end{equation}
We extend $p_t^j$ by zero outside $D_t^j$. It is a probability
vector on that set, so
$0\le\Delta_t^j=\sum_xp_t^j(x)^2-1/r_j\le1$.

For $0\le t<T$ and a feasible direction prefix write
$K_t(v)=\Pr\{v_{t+1}=v\mid v_1,\ldots,v_t\}$. Conditional on the
complete schedule, the law of the \emph{joint} permutation path
through time $t$ depends only on $v_1,\ldots,v_t$. The independence
of $\Xi$ therefore gives
\begin{equation}\label{eq:marginal-direction-factorization}
 \Pr\{Z_t(c_j)=x,\ v_{t+1}=v\mid\mathcal F_t^j\}
 =p_t^j(x)K_t(v).
\end{equation}
Indeed, the joint law of the old hidden position and observed path
prefix is the same for every completion of a fixed direction
prefix. Conditioning on that observed prefix leaves the
conditional law of the future directions unchanged. This argument
uses only path prefixes; it does not condition on future observed
paths or on further auxiliary data.

\begin{lemma}[Online averaging and cross-hyperplane contraction]
\label{lem:marginal-averaging}\label{shear:lem:energy}
In the preceding setup, for $0\le t<T$ and $1\le j\le q$,
\begin{equation}\label{eq:marginal-exact-decrement}
 \Delta_{t+1}^j=\Delta_t^j-\frac12
 \sum_{\substack{\{x,y\}\subseteq D_t^j\\y=x+v_{t+1}}}
       \bigl(p_t^j(x)-p_t^j(y)\bigr)^2
 \quad\text{almost surely},
\end{equation}
where each unordered matching edge is counted once. If $K_t$ is
uniform on the complement of an $\mathcal F_t^j$-measurable
hyperplane $J_t$, then
\begin{equation}\label{eq:marginal-cross-energy}
 \E[\Delta_t^j-\Delta_{t+1}^j\mid\mathcal F_t^j]
 \ge
 \frac{\min(|D_t^j\cap J_t|,|D_t^j\setminus J_t|)}{n}\Delta_t^j.
\end{equation}
More generally, suppose $K_t$ is uniform outside an
$\mathcal F_t^j$-measurable proper subspace $W_t$, and let $J_t$
be an $\mathcal F_t^j$-measurable hyperplane containing $W_t$.
Then the same right side with denominator $2n$ is a valid lower
bound.
\end{lemma}
\begin{proof}
Fix $j$ and omit its superscript. By
\eqref{eq:marginal-direction-factorization}, revealing $v_{t+1}$
does not update the old hidden-card law $p_t$. Given
$\mathcal F_t$ and this direction, the matching edges incident to
observed cards are known. Their next positions reveal exactly the
fresh coins on those edges. Those coins are independent of the old
hidden position, so their revelation also leaves $p_t$ unchanged;
all other edge coins remain independent and fair.

An edge containing two positions of $D_t$ averages their two
masses. An edge containing exactly one position of $D_t$ transports
its mass to the unique position on that edge unoccupied by observed
cards after the switch. Edges with no available position carry no
hidden mass. These rules give the conditional vector at the exact
field $\mathcal F_{t+1}$ in \eqref{eq:marginal-filtration}. They
also transport the uniform vector on $D_t$ to the uniform vector
on $D_{t+1}$. Transport preserves squared error, while averaging
two masses reduces it by
$\frac12(p_t(x)-p_t(y))^2$. This proves the exact identity.

For the expectation, split $D_t=D^0\cup D^1$ across $J_t$ and put
$z_x=p_t(x)-1/r_j$, $u=|D^0|$, and $w=|D^1|$. Since $\sum z_x=0$,
\begin{equation}\label{shear:eq:bipartite-form}
 \sum_{x\in D^0,y\in D^1}(z_x-z_y)^2
 =w\sum_{D^0}z_x^2+u\sum_{D^1}z_y^2
     +2\left(\sum_{D^0}z_x\right)^2
 \ge\min(u,w)\sum_{D_t}z_x^2.
\end{equation}
Every cross pair has matching probability $2/n$ in the hyperplane
case. Combining this with the factor $1/2$ in the exact decrement
proves \eqref{eq:marginal-cross-energy}. In the proper-subspace
case, every cross difference lies outside $W_t$ and has probability
$1/(n-|W_t|)\ge1/n$; discard the other eligible pairs to obtain
the denominator $2n$. The calculation covers an empty side.
The algebraic identity \eqref{shear:eq:bipartite-form} holds for
any real centered vector; a later use for another vector must also
identify its averaging rule.
\end{proof}

The loss bound and a balance estimate will be used at different
levels of conditioning. Let $0\le\gamma\le1$. If the conditional expected loss
$\E[\Delta_t^j-\Delta_{t+1}^j\mid\mathcal F_t^j]$ is at least
$\gamma\Delta_t^j$ off an event $B_t\in\mathcal F_t^j$ and is
always nonnegative, then $\Delta_t^j\le1$ gives
\begin{equation}\label{eq:marginal-balance-recurrence}
 \E\Delta_{t+1}^j
 \le(1-\gamma)\E\Delta_t^j+\gamma\Pr(B_t).
\end{equation}
Thus an unconditional bound on $\Pr(B_t)$ suffices; no bound
conditional on the observed paths is asserted.

\subsection{From hidden cards to an ordinary tuple law}

Let $\pi_q$ be uniform on ordered distinct $q$-tuples of positions,
and let $U_D$ be uniform on a nonempty set $D$. For any tuple law
$\nu=\mathcal L(Y_1,\ldots,Y_q)$,
\begin{equation}\label{eq:marginal-sequential-tv}
 \|\nu-\pi_q\|_{\TV}
 \le\sum_{j=1}^q\E_\nu
 \left\|
 \mathcal L(Y_j\mid Y_1,\ldots,Y_{j-1})
 -U_{V\setminus\{Y_1,\ldots,Y_{j-1}\}}
 \right\|_{\TV}.
\end{equation}
This sequential comparison appears in
\cite[Appendix~A, Lemma~2]{mrs}; an earlier distinct-tuple form is
\cite[Section~5.3, Lemma~12]{morris-exclusion2006}.
For completeness, at the $j$th step insert the law that first
draws the actual $(j-1)$-coordinate prefix of $\nu$ and then
appends a uniform available position. Its distance from the
$j$-coordinate marginal of $\nu$ is the $j$th summand above.
Applying the same extension kernel to $\pi_{j-1}$ gives $\pi_j$
and contracts total variation. The triangle inequality and
induction prove \eqref{eq:marginal-sequential-tv}. Every
expectation is under the actual prefix law.

\begin{lemma}[Terminal data and an ordinary tuple law]
\label{shear:lem:tuple}
In the preceding setup let $\Lambda$ be terminal data that
determine the complete direction schedule, and suppose $\Xi$ is
independent of $(\Lambda,(Z_t)_{t=0}^T)$. For every
$1\le j\le q$ and every $x\in V$, assume
that its online terminal vector is also the conditional hidden
law after these data and the complete predecessor paths are known:
\begin{equation}\label{eq:marginal-terminal-identification}
 \Pr\{Z_T(c_j)=x\mid
       \Lambda,\Xi,\,
       Z_s(c_i):0\le s\le T,\ i<j\}
 =p_T^j(x)\quad\text{almost surely}.
\end{equation}
Then
\begin{equation}\label{shear:eq:tuple-bound}
 \E_{\Lambda,\Xi}
 \|\mathcal L(Z_T(C)\mid\Lambda,\Xi)-\pi_q\|_{\TV}
 \le\frac12\sum_{j=1}^q
       \sqrt{(n-j+1)\E\Delta_T^j}.
\end{equation}
All expectations are under the actual joint law of the data, list,
and process. If $\E\Delta_T^j\le a_T$ for every $j$, the right
side is at most $q\sqrt n\,\sqrt{a_T}/2$. For $q=0$ both sides are
zero.
\end{lemma}
\begin{proof}
For fixed $\Lambda,\Xi$, apply
\eqref{eq:marginal-sequential-tv}. At index $j$ use the coarser
field
\[
 \mathcal H_j
 =\sigma(\Lambda,\Xi,Z_T(c_1),\ldots,Z_T(c_{j-1})).
\]
The tower property and \eqref{eq:marginal-terminal-identification}
give
$\Pr\{Z_T(c_j)=x\mid\mathcal H_j\}
 =\E[p_T^j(x)\mid\mathcal H_j]$.
The set $D_T^j$ is $\mathcal H_j$-measurable, so Jensen's
inequality for the $\ell^1$ norm compares both vectors to the
same uniform law $U_{D_T^j}$. Cauchy--Schwarz gives
\[
 \|p_T^j-U_{D_T^j}\|_{\TV}
 \le\frac12\sqrt{(n-j+1)\Delta_T^j}.
\]
Consequently, almost surely in the terminal data and initial data,
\begin{equation}\label{eq:marginal-conditional-tuple}
 \|\mathcal L(Z_T(C)\mid\Lambda,\Xi)-\pi_q\|_{\TV}
 \le\frac12\sum_{j=1}^q
 \E\!\left[\sqrt{(n-j+1)\Delta_T^j}\,\middle|\,\Lambda,\Xi\right].
\end{equation}
Averaging and applying Cauchy--Schwarz once more proves the lemma.
\end{proof}

When $\Lambda=\mathbf v$, the field on the left of
\eqref{eq:marginal-terminal-identification} is exactly
$\mathcal F_T^j$, so that hypothesis is automatic. If $\Lambda$
contains further shear or functional data, the identification
requires an additional path-law argument. We will prove that
argument where those data are introduced; their disclosure is
separate from the chronological averaging above.

\subsection{A sampled omitted block}

For the local full-group proof, set
\begin{equation}\label{eq:marginal-constants}
 k=64,\qquad m=\frac{n}{k},\qquad
 \alpha=\frac1{4k},\qquad T=(1+16k)d=1025d,
\end{equation}
and assume $d\ge6$. Put
\begin{equation}\label{eq:marginal-error}
 b_n=2\left(\frac{2}{3^{3/4}}\right)^m,\qquad
 \varepsilon_n=\frac{n^{3/2}}2\sqrt{e^{-4d}+b_n}.
\end{equation}
For a labeled partition $\mathcal M=(M_1,\ldots,M_k)$ of $V$
into $m$-element blocks, let $\delta_a(\mathcal M)$ be the
distance of the time-$T$ images of $V\setminus M_a$ under the
physical shuffle from the uniform distinct tuple law. Any fixed
ordering of these labels gives the same distance.

\begin{proposition}[One partition with nearly complete marginals]
\label{prop:marginals}
There is a deterministic labeled partition
$\mathcal M=(M_1,\ldots,M_k)$ into blocks of size $m$ for which
\begin{equation}\label{eq:marginal-partition}
 \sum_{a=1}^k\delta_a(\mathcal M)
 \le k\varepsilon_n=o(1)\qquad(d\longrightarrow\infty).
\end{equation}
\end{proposition}

\begin{proof}
Use $v_i=e_i$ for $1\le i\le d$. For $d\le t<T$ put
\begin{equation}\label{eq:marginal-hyperplane}
 J_t=\Span\{v_{t-d+2},\ldots,v_t\}
\end{equation}
and choose $v_{t+1}$ uniformly in $V\setminus J_t$, with fresh
direction randomness. Every $d$ successive directions form a
basis. Conditional on the complete schedule, use independent fair
switches to form $Z$. Independently choose $\mathcal M$ uniformly
among labeled equal-size partitions. For a fixed $a$, let $\Xi$
be this partition and let $C$ be any deterministic ordering of
$V\setminus M_a$, with $q=n-m$. This is the online setup above.

For each hidden index $j$, the available set contains the positions
of all omitted labels:
\begin{equation}\label{eq:marginal-permanent-holes}
 Z_t(M_a)\subseteq D_t^j.
\end{equation}
Write $D_t^{j,0}=D_t^j\cap J_t$ and
$D_t^{j,1}=D_t^j\setminus J_t$. If
$J_t^0=J_t$ and $J_t^1=V\setminus J_t$, then
\[
 \left\{\min_i|D_t^j\cap J_t^i|<m/4\right\}
 \subseteq
 \left\{\min_i|Z_t(M_a)\cap J_t^i|<m/4\right\}.
\]
Fix the direction and switch history through time $t$, but not
$\mathcal M$. This fixes $Z_t$ and $J_t$, while $Z_t(M_a)$ is a
uniform $m$-subset of $V$.

For a uniform $m$-subset and a fixed half of $V$, let $X$ be its
count in that half. View the subset as an ordered sample without
replacement. A specified set of $\ell$ sampled indices all lie
in the half with probability at most $2^{-\ell}$. If $I_i$ are
the membership indicators, expansion gives
\[
 \E3^X=\E\prod_{i=1}^m(1+2I_i)
 \le\sum_{A\subseteq\{1,\ldots,m\}}2^{|A|}2^{-|A|}=2^m.
\]
If one half contains fewer than $m/4$ sampled points, the other
contains more than $3m/4$. Markov's inequality and a union bound
therefore yield, for $d\le t<T$,
\begin{equation}\label{eq:marginal-balance}
 \Pr\{\min(|D_t^{j,0}|,|D_t^{j,1}|)<m/4\}
 \le2\frac{2^m}{3^{3m/4}}=b_n.
\end{equation}
This bound averages the independent partition and the shuffle; it
is not a concentration assertion conditional on the observed paths.

Off the event in \eqref{eq:marginal-balance}, the conditional loss
in \eqref{eq:marginal-cross-energy} is at least
$\alpha\Delta_t^j$. Its loss is always nonnegative. Applying
\eqref{eq:marginal-balance-recurrence} gives
\begin{equation}\label{eq:marginal-recursion}
 \E\Delta_{t+1}^j
 \le(1-\alpha)\E\Delta_t^j+\alpha b_n.
\end{equation}
Starting with $\Delta_d^j\le1$ and iterating for $T-d=16kd$
steps, we obtain, uniformly in $a,j$,
\begin{equation}\label{eq:marginal-terminal-energy}
 \E\Delta_T^j\le(1-\alpha)^{T-d}+b_n
 \le e^{-4d}+b_n.
\end{equation}

Let $\widetilde\delta_a(\mathbf v,\mathcal M)$ be the auxiliary
tuple distance conditional on the complete schedule and partition.
Use Lemma~\ref{shear:lem:tuple} with $\Lambda=\mathbf v$.
Its terminal identification is automatic, since
$\mathcal F_T^j$ already contains the complete schedule and
predecessor paths. Equations
\eqref{shear:eq:tuple-bound} and
\eqref{eq:marginal-terminal-energy} give
\begin{equation}\label{eq:marginal-averaged-tuple}
 \E_{\mathbf v,\mathcal M}
 \widetilde\delta_a(\mathbf v,\mathcal M)
 \le\frac{n^{3/2}}2\sqrt{e^{-4d}+b_n}
 =\varepsilon_n.
\end{equation}
The schedule remains inside the distance on the left.

For every fixed schedule and partition, choose
$\boldsymbol\iota=(1,\ldots,d)$ in Lemma~\ref{lem:frames}.
The first directions are $v_j=e_j$, so $L_0=I$. Corollary
\ref{cor:frames:marginals}, followed by the terminal rotation,
therefore gives the same-list identity
\begin{equation}\label{eq:marginal-frame-transfer}
 \widetilde\delta_a(\mathbf v,\mathcal M)=\delta_a(\mathcal M).
\end{equation}
The frame maps enter only this fixed-schedule endpoint comparison.
Summing \eqref{eq:marginal-averaged-tuple} over $a$ now gives
$\E_{\mathcal M}\sum_a\delta_a(\mathcal M)\le k\varepsilon_n$.
At least one deterministic partition attains this average bound.
Finally,
$n^{3/2}e^{-2d}=\exp((\tfrac32\log2-2)d)\to0$, while
$2/3^{3/4}<1$ makes $b_n$ exponentially small in $n$.
Thus $\varepsilon_n\to0$.
\end{proof}
